% eksperymentalne tłumaczenie na włoski

%

\subsection{Okręgi Apoloniusza i Soddy'ego} % lang-pl

\subsection{Cerchi di Apollonio e di Soddy} % lang-it

% https://en.wikipedia.org/wiki/Problem_of_Apollonius
\begin{problem}
    \label{problem_apoloniusza}%
    Na płaszczyźnie dane są trzy okręgi. % lang-pl
    Skonstruować czwarty, który jest styczny do każdego z~pozostałych. % lang-pl
    
    \index{okrąg!Apoloniusza}% % lang-pl
    Nel piano sono dati tre cerchi. % lang-it
    Costruire un quarto che sia tangente a ciascuno degli altri. % lang-it
    \index{cerchio!di Apollonio}% % lang-it
\end{problem}

Dyskusję na temat problemu znajdziemy u Guzickiego \cite[s. 444-461]{guzicki_2021}. % lang-pl
Eves \cite[s. 164]{eves1_1972} rozważa szczególny przypadek trzech współpękowych okręgów. % lang-pl
W ogólności okręgi mogą mieć zerowy promień (i skurczyć się do punktu) lub nieskończony promień (i wydłużyć się do prostej). % lang-pl
Zależnie od wyboru kształtów, mamy do rozwiązania dziesięć różnych zadań: PPP, PPL, PLL, LLL, PPC, PLC, LLC, PCC, LCC, CCC (klasyczne zadanie Apoloniusza). % lang-pl
Tutaj P oznacza punkt, L prostą, zaś C okrąg, od pierwszych liter stosownych słów po angielsku.% % lang-pl
Una discussione del problema si trova in Guzicki \cite[s. 444-461]{guzicki_2021}. % lang-it
Eves \cite[s. 164]{eves1_1972} considera il caso particolare di tre cerchi coassiali. % lang-it
In generale i cerchi possono avere raggio nullo (riducendosi a un punto) oppure raggio infinito (diventando una retta). % lang-it
A seconda della scelta delle figure, si ottengono dieci problemi diversi: PPP, PPL, PLL, LLL, PPC, PLC, LLC, PCC, LCC, CCC (il problema classico di Apollonio). % lang-it
Qui P indica un punto, L una retta e C un cerchio, dalle iniziali dei termini inglesi corrispondenti.% % lang-it

Z zadaniem wiąże się dużo późniejszy wynik. % lang-pl
W dwóch listach z 1643 roku do czeskiej księżniczki Elżbiety z Palatynatu Kartezjusz przedstawi uproszczoną wersję, w której trzy okręgi są parami styczne i znajdzie (bez dowodu) relację wiążącą promienie okręgów.% % lang-pl
\index[persons]{van Pallandt, Elżbieta (z Palatynatu)}% % lang-pl
Al problema è legato un risultato molto più tardo. % lang-it
In due lettere del 1643 alla principessa Elisabetta di Boemia, Cartesio presentò una versione semplificata in cui tre cerchi sono a due a due tangenti e trovò (senza dimostrazione) una relazione tra i loro raggi.% % lang-it
\index[persons]{di Boemia, Elisabetta (Principessa Palatina)}% % lang-it

\begin{theorem}[Kartezjusza] % lang-pl
Dane są trzy okręgi, parami styczne o promieniach $r_1, r_2, r_3$ oraz czwarty styczny do każdego z nich\footnote{istnieją dwa takie, jeden z~nich zazwyczaj styczny wewnętrznie, a w szczególnym przypadku może okazać się prostą.}. % lang-pl
Definiujemy krzywiznę ze znakiem jako $\pm 1 / r$, gdzie znak $-$ odpowiada okręgowi stycznemu zewnętrznie. % lang-pl
Wtedy% % lang-pl
\begin{theorem}[di Cartesio] % lang-it
Sono dati tre cerchi, a due a due tangenti, con raggi $r_1, r_2, r_3$ e un quarto tangente a ciascuno di essi\footnote{ne esistono due, uno dei quali è di solito tangente internamente e, in un caso particolare, può degenerare in una retta.}. % lang-it
Definiamo la curvatura con segno come $\pm 1 / r$, dove il segno $-$ corrisponde a una tangenza esterna. % lang-it
Allora% % lang-it
\begin{equation}
    \left(\pm \frac{1}{r_1} \pm  \frac{1}{r_2} \pm  \frac{1}{r_3} \pm  \frac{1}{r_4}\right)^2 = 2 \left(\frac{1}{r_1^2} + \frac{1}{r_2^2} + \frac{1}{r_3^2} + \frac{1}{r_4^2}\right),
\end{equation}
gdzie $r_4$ oznacza promień czwartego okręgu.% % lang-pl
dove $r_4$ indica il raggio del quarto cerchio.% % lang-it
\end{theorem}

Jeżeli jeden z okręgów (dajmy na to ten o promieniu $r_1$) staje się prostą, to jego zerowa krzywizna znika z równania. % lang-pl
Jeżeli dwa okręgi stają się dwiema prostymi, to pozostałem dwa muszą być do siebie przystające, co należy uznać za zdegenerowany przypadek.% % lang-pl
Se uno dei cerchi (ad esempio quello di raggio $r_1$) diventa una retta, la sua curvatura nulla scompare dall'equazione. % lang-it
Se due cerchi diventano due rette, gli altri due devono essere congruenti, il che va considerato un caso degenere.% % lang-it

Twierdzenie Kartezjusza odkryją na nowo Yamaji Nuchizumi (1751), Jakob Steiner (\emph{Fortsetzung der geometrischen Betrachtungen} z 1826), Philip Beecroft (\emph{Properties of circles in mutual contact} z 1842) i~wreszcie Frederick Soddy (\emph{The Kiss Precise} z 1936 jak opisane u Coxetera \cite[s. 29-31]{coxeter_1967}).% % lang-pl
Il teorema di Cartesio fu riscoperto da Yamaji Nuchizumi (1751), Jakob Steiner (\emph{Fortsetzung der geometrischen Betrachtungen}, 1826), Philip Beecroft (\emph{Properties of circles in mutual contact}, 1842) e infine Frederick Soddy (\emph{The Kiss Precise}, 1936, come descritto in Coxeter \cite[s. 29-31]{coxeter_1967}).% % lang-it
\index[persons]{Nuchizumi, Yamaji}%
\index[persons]{Steiner, Jakob}%
\index[persons]{Beecroft, Philip}%
\index[persons]{Soddy, Frederick}%

Ten ostatni poda uogólnienie do sfer, a rok później Thorold Gosset przeskoczy do przestrzeni dowolnego wymiaru ($\mathbb R^n$), z $n$ w miejsce czynnika $2$.% % lang-pl
Quest'ultimo fornì una generalizzazione alle sfere, e un anno dopo Thorold Gosset passò allo spazio di dimensione arbitraria ($\mathbb R^n$), con $n$ al posto del fattore $2$.% % lang-it
\index[persons]{Gosset, Thorold}%

Dowody korzystają z inwersji, wzoru Herona, wyznacznika Cayleya-Mengera, albo łańcuchów Pappusa i twierdzenia Vivianiego. % lang-pl
Patrz też $\Delta_{10}^6$.% % lang-pl
Le dimostrazioni utilizzano inversioni, la formula di Erone, il determinante di Cayley--Menger, oppure le catene di Pappo e il teorema di Viviani.% % lang-it
Si veda anche $\Delta_{10}^6$. % lang-it

\begin{example}
Dane są trzy styczne zewnętrznie okręgi o promieniach $r_1 = r_2 = r_3 = \sqrt{3}$. % lang-pl
Czwarty okrąg styczny do nich wszystkich ma wtedy krzywiznę $\sqrt 3 \pm 2$.% % lang-pl
Sono dati tre cerchi tangenti esternamente con raggi $r_1 = r_2 = r_3 = \sqrt{3}$. % lang-it
Il quarto cerchio tangente a tutti ha allora curvatura $\sqrt 3 \pm 2$.% % lang-it
\end{example}

Jeszcze nie skończyliśmy z Soddym. % lang-pl
W wierzchołkach trójkąta o obwodzie $2p = a + b + c$ wstawiamy okręgi o promieniach $p - a$, $p - b$, $p - c$. % lang-pl
Wtedy czwarty okrąg (zwany okręgiem Soddy'ego) z twierdzenia Kartezjusza ma promień% % lang-pl
Non abbiamo ancora finito con Soddy. % lang-it
Nei vertici di un triangolo di perimetro $2p = a + b + c$ inseriamo cerchi di raggi $p - a$, $p - b$, $p - c$. % lang-it
Allora il quarto cerchio (detto cerchio di Soddy) dal teorema di Cartesio ha raggio% % lang-it
\begin{equation}
    \frac{4R + r \pm 2p}{S},
\end{equation}
gdzie $R, r$ to promienie okręgu opisanego i wpisanego, zaś $S$ to pole powierzchni. % lang-pl
Napisze o tym Coxeter \cite[s. 29-31]{coxeter_1967}.% % lang-pl
dove $R, r$ sono i raggi del cerchio circoscritto e inscritto, e $S$ è l'area. % lang-it
Ne scrive Coxeter \cite[s. 29-31]{coxeter_1967}.% % lang-it

\begin{definition}[prosta Soddy'ego] % lang-pl
    Prostą przechodzącą przez środki okręgów Soddy'ego nazywamy prostą Soddy'ego. % lang-pl
\begin{definition}[retta di Soddy] % lang-it
    
    La retta che passa per i centri dei cerchi di Soddy si chiama retta di Soddy. % lang-it
\end{definition}

\begin{proposition}
Prosta Soddy'ego przechodzi przez środek okręgu wpisanego, punkt Gergonne'a i punkt Eppsteina (punkt, gdzie przecinają się trzy proste przechodzące przez trzy pary spośród sześciu punktów styczności czterech okręgów...)% % lang-pl
\index{punkt!Eppsteina}% % lang-pl
La retta di Soddy passa per il centro del cerchio inscritto, il punto di Gergonne e il punto di Eppstein (il punto di intersezione di tre rette che passano per tre coppie dei sei punti di tangenza dei quattro cerchi...).% % lang-it
\index{punto!di Eppstein}% % lang-it
\end{proposition}

Prosta Soddy'ego ma pewien związek z punktem de Lonchampsa, którego nie opiszemy.% % lang-pl
\index{punkt!de Longchampsa}% % lang-pl
La retta di Soddy ha una certa relazione con il punto di de Longchamps, che non descriveremo.% % lang-it
\index{punto!di de Longchamps}% % lang-it

%